% Options for packages loaded elsewhere
\PassOptionsToPackage{unicode}{hyperref}
\PassOptionsToPackage{hyphens}{url}
%
\documentclass[
]{article}
\usepackage{amsmath,amssymb}
\usepackage{lmodern}
\usepackage{iftex}
\ifPDFTeX
  \usepackage[T1]{fontenc}
  \usepackage[utf8]{inputenc}
  \usepackage{textcomp} % provide euro and other symbols
\else % if luatex or xetex
  \usepackage{unicode-math}
  \defaultfontfeatures{Scale=MatchLowercase}
  \defaultfontfeatures[\rmfamily]{Ligatures=TeX,Scale=1}
\fi
% Use upquote if available, for straight quotes in verbatim environments
\IfFileExists{upquote.sty}{\usepackage{upquote}}{}
\IfFileExists{microtype.sty}{% use microtype if available
  \usepackage[]{microtype}
  \UseMicrotypeSet[protrusion]{basicmath} % disable protrusion for tt fonts
}{}
\makeatletter
\@ifundefined{KOMAClassName}{% if non-KOMA class
  \IfFileExists{parskip.sty}{%
    \usepackage{parskip}
  }{% else
    \setlength{\parindent}{0pt}
    \setlength{\parskip}{6pt plus 2pt minus 1pt}}
}{% if KOMA class
  \KOMAoptions{parskip=half}}
\makeatother
\usepackage{xcolor}
\setlength{\emergencystretch}{3em} % prevent overfull lines
\providecommand{\tightlist}{%
  \setlength{\itemsep}{0pt}\setlength{\parskip}{0pt}}
\setcounter{secnumdepth}{-\maxdimen} % remove section numbering
\ifLuaTeX
  \usepackage{selnolig}  % disable illegal ligatures
\fi
\IfFileExists{bookmark.sty}{\usepackage{bookmark}}{\usepackage{hyperref}}
\IfFileExists{xurl.sty}{\usepackage{xurl}}{} % add URL line breaks if available
\urlstyle{same} % disable monospaced font for URLs
\hypersetup{
  hidelinks,
  pdfcreator={LaTeX via pandoc}}

\author{Dietmar Gerald Schrausser}
\date{09.05.2026}


\begin{document}

\hypertarget{foliumblat-iir}{%
\section{Folium/Blat IIr}\label{foliumblat-iir}}

A \emph{contemporary} English translation is presented from comparing
the existing Latin, Early New High German (ENHG) and English
(translation) texts for better comprehensibility. Especially since the
ENHG text is difficult to understand compared to modern High German, as
a \emph{profound} knowledge of German (as a mother tongue) as well as
already aknowledge of various \emph{dialects} is required.

\hypertarget{without-heading.}{%
\subsection{Without heading.}\label{without-heading.}}

\begin{quote}
IN principio creauit deus celum et terra{[}m{]}. Terra autem erat inanis
{[}et{]} vacua. {[}et{]} tenebrae era{[}n{]}t sup{[}er{]} faciem abissi.
et sp{[}irit{]}us d{[}omi{]}ni ferebat{[}ur{]} sup{[}er{]} aq{[}ua{]}s.
\end{quote}

\textbf{Moses}, the divine prophet and historian who lived nearly 700
years before the Trojan War, teaches that God, the planner and creator
of all things, first created the heaven and established it as the seat
of the Creator God himself in order to accomplish this work. Then he
founded the earth and subordinated it to the heaven. He shrouded the
earth in darkness, meaning it is without its own enlightment\footnote{`dan{[}n{]}
  sie begreüft durch sich selbs nichtzit'.} unless it receives light
from heaven. There he instituted eternal light, the heavenly spirits and
eternal life; to the earth, however, he gave darkness, earthly spirits
and death.\footnote{from Lactantius
  (\href{https://books.google.com/books?id=qGsnghsi3uUC}{1497}, Book 2,
  Chapter 9/10).}

By describing a divine \emph{creation}, Moses exposes the three
errors\footnote{Refers to the classical philosophical views on the
  creation and eternity of the universe that contradicted Genesis.} of
\textbf{Plato}, \textbf{Aristotle} and \textbf{Epicurus}. Plato believed
that (1) God, (2) the \emph{Ideas}\footnote{`ydeas' and `die vorpildnus
  oder gestaltnus seiner geschöpff', ideas (Ydeas) are the eternal,
  perfect archetypes in Platonic thought.} and Yle had existed from
eternity and (3) that the world was fundamentally made from this Yle.
The Greeks used the term \emph{Yle} to refer to the first formless mass
from which \emph{everything was created}; the \emph{visible} things
formed from harmonizied elements. Or, the finest dust\footnote{`de
  athomis' and `de{[}m{]} aller dynnisten staub'.} that glitters in
sunlight, made from matter and form.\footnote{attributed to the `Summa
  super Sententias' (Petrus,
  \href{https://doi.org/10.3931/e-rara-18241}{1484}, Book II, Section
  1), it deals with Plato's view of eternal, unformed matter,
  Aristotle's claim of an eternal world and Epicurus's theory of atomic
  collision.}

(ad 3) Nevertheless\footnote{`t{[}ame{]}n'.}, God created the world
without any available and pre-prepared material; for \emph{he} is to be
regarded as the most cunning\footnote{`prude{[}n{]}tissim{[}us{]}'.} and
wisest Creator before he accomplished the work of the world. In himself
was the source of perfect and accomplished goodness, so that\footnote{`vt'.}
goodness might flow forth from it like a
\emph{stream}\footnote{`tam{[}que{]}m riu{[}us{]}' and `als ein pach'.}.\footnote{It
  serves as a Christian-Neoplatonic argument for creation from nothing
  (\emph{creatio ex \textbf{nihilo}}) and contrasts it with Greek
  philosophical views. It emphasizes that \emph{God} acted not
  accidentally, but \emph{consciously}, `excogitandum' (c.f.
  Trismegistu,
  \href{https://gdz.sub.uni-goettingen.de/id/PPN858183994}{1471}).}

(ad 2) Of all beings, he primordially created the
\emph{angels}\footnote{The intellectual world, s. fol.~1r `quem theologi
  angelicum: philosophi: autem intellectualem vocant'.} from that which
is \emph{nothing} (\textbf{\emph{n{[}{]}}}).\footnote{`ex eo q{[}uo{]}d
  n{[}{]} e{[}st{]}' and `auß de{[}n{]} das nicht ist'.} Thus, \emph{he}
is powerful through \emph{eternity} and, through the strength of this
immense power (which is limitless and immeasurable\footnote{`q{[}uam{]}
  fine ac mo{[}do{]} caret' and `die des ends vn{[}d{]} der maß
  mangelt'.}), \emph{equal} to the life of the
Creator\footnote{`sic{[}ut{]} vita facturis' vs.~`als das lebe{[}n{]}
  des Schöpffers'.}.\footnote{discusses the Christian theological
  concept of creation from nothing (\emph{creatio ex nihilo}).}
Therefore, it may seem astonishing that, prior to the creation of the
world, raw matter\footnote{`materia{[}m{]} de q{[}ua{]} faceret' and
  `ein materi daraus er machet'.} from \emph{nothing} had to be
prepared.\footnote{c.f. Lactantius
  (\href{https://books.google.com/books?id=qGsnghsi3uUC}{1497}, Book II
  8.8), he argued against ancient Greek philosophies claiming God
  created the world from \emph{pre-existing} matter.} This is probably
how the Saracens understood\footnote{from the `Etymologiae' (De Seville,
  \href{https://books.google.com/books?id=wb8Z_vlSyJwC}{1802}, Book
  VIII, Chapter xi, \emph{De diis gentium}), written in the \emph{early}
  7th c., around 610 AD. In this case, it can be assumed that the
  reference is to the \emph{Saracens} from the 4th c., nomadic tribes of
  the Syrian-Arabian desert, and therefore there is no reference to
  \emph{Islam}. Rather, it seems to refer to the \emph{pivotal} events
  around 378 AD, in connection with the devastating defeat of the
  Western Roman Empire against the Goths at Adrianople and the
  involvement of the \emph{Saracens} under \emph{Mavia}
  (\emph{Mawiyya}), a female Arab ruler who had recently led a revolt
  against the Roman Empire in southern Syria and subsequently
  \emph{switched sides}. Interestingly, the story of the \emph{Saint}
  \textbf{\emph{Moses}} \emph{the Bishop of the Arabs} also falls into
  this context - not to be confused with the similarly named
  \textbf{\emph{Moses}} \emph{the Black}, a converted criminal who
  operated in the same region at the same time (also known as
  \emph{Moses the Strong}, \emph{Moses the Abyssinian}, \emph{Moses the
  Indian}, \emph{Abba Moses the Robber}, \emph{Moses Murin} or
  \emph{Moses the Ethiopian} etc.).} it when they said that the angels
had been led by God from \emph{darkness} into the light and filled with
eternal joy. However, some of them forgot their divine
character\footnote{`indolis' and `einpildung'.}, changed from good to
evil\footnote{`ex bono p{[}er{]} se malu{[}m{]} effectu{[}m{]}' and
  `vo{[}m{]} gůtten zum ubel getretten'.} through their own \emph{own}
reversal\footnote{`verkerung', perversion.} and became devils (what we
call criminals\footnote{`greci diabolu{[}m{]} appelant: nos
  crimi{[}n{]}atore{[}m{]} vocam{[}us{]}'.}).\footnote{from the
  `Heptaplus' (Mirandola \& Mirandola,
  \href{https://doi.org/10.3931/e-rara-61217}{1601}), Book 1, Chapter
  1).}

(ad 1) The earth was void, that is (as \textbf{Jerome} and \textbf{the
Seventy}\footnote{`septinge{[}n{]}ta', 700 vs.~`die .lxx.
  auslege{[}r{]}'.} suggest) invisible and unjoined\footnote{`i{[}n{]}co{[}m{]}posita'
  and `vnzesamen gefügt', refers to the Latin translations of the Hebrew
  `tohu wa-bohu' (formless and empty) from Genesis 1:2 (c.f. Jiménez de
  Cisneros, \href{https://doi.org/10.3931/e-rara-46695}{1517}).} and
because of this confused\footnote{`{[}con{]}fusio{[}n{]}e' and
  `zestrewlichkeit'.} state called \emph{Abyss}, which the Greeks call
\emph{Chaos}. This is\footnote{`i{[}d est{]}'.} \emph{three-dimensional}
matter\footnote{`materia{[}m{]} trino dime{[}n{]}su' and `materi mit
  driueltiger ermessung'.}, extended to the greatest depths\footnote{`i{[}n{]}
  altissimas profunditates extensam'.}.

Ovid also remembers this:

\begin{quote}
Ante mare {[}et{]} terras {[}et{]} q{[}uo{]}d tegit o{[}mn{]}ia
celu{[}m{]}. Un{[}us{]} erat toto nature vult{[}us{]} in orbe.
Que{[}m{]} dixere chaos rudis indigestaq{[}ue{]} moles. Nec
q{[}ui{]}cq{[}uam{]} nisi po{[}n{]}dus iners: {[}con{]}gestaq{[}ue{]}
eode{[}m{]}. No{[}n{]} bene iu{[}n{]}ctaru{[}m{]} discordia femina
reru{[}m{]}. Nullus ad huc mu{[}n{]}do p{[}re{]}bebat lumina titan.
\end{quote}

Thus, the Spirit of the Lord (as an instrument of divine art) floated
over the waters, like the will of a master builder who organizes all
that needs to be done. Through this, God's works are perfect and
creation is expressed in the \emph{number six}\footnote{`senario numero'
  and `in sechser zall'.}, whose components are 1, 2 and 3, which rise
to form a triangle\footnote{`que in trigonu{[}m{]} surga{[}n{]}t'.}. In
his account of the six days, Moses dedicates the first day (i) to
creation, the second and third days (ii) to order and design and (iii)
the remaining days to the adornment of the world.\footnote{from the
  `Heptaplus', (Mirandola \& Mirandola,
  \href{https://doi.org/10.3931/e-rara-61217}{1601}), in classical and
  medieval number theory, 6 is the first \emph{perfect number} because
  it is equal to the sum of its proper divisors (1+2+3=6).}

\hypertarget{references}{%
\subsection{References}\label{references}}

De Seville, I. (1802). \emph{S. Isidori Hispalensis Episcopi Hispaniarum
Doctoris Opera Omnia}. Edited by Arevalo, F. Romae: Apvd Antonivm
Fvlgonivm. \url{https://books.google.com/books?id=wb8Z_vlSyJwC}

Jiménez de Cisneros, F. (1517). \emph{Biblia Polyglotta Complutensis}.
Complutum: Arnaldo Guillén de Brocar.
\url{https://doi.org/10.3931/e-rara-46695}

Lactantius, L.C.F. (1497). \emph{Divinae Institutiones}. Venetiis: Simon
Bevilaqua. \url{https://books.google.com/books?id=qGsnghsi3uUC}

Mirandola, G., \& Mirandola, G. F. (1601). \emph{Ioannis Pici,
Mirandulae \ldots{} opera quae extant omnia \ldots{}} Basileae: per
Sebastianum Henricpetri. \url{https://doi.org/10.3931/e-rara-61217}

Petrus. (1484). \emph{Sententiarum Libri IV}. Basel.
\url{https://doi.org/10.3931/e-rara-18241}

Trismegistu, H. (1471). \emph{Mercurii Trismegisti Liber de Potestate Et
Sapientia Dei e Graeco in Latinum Traductus a Marsilio Ficino Pimander
Incipit}. Edited by Ficinus, M. Treviso.
\url{https://gdz.sub.uni-goettingen.de/id/PPN858183994}

\end{document}
